\section{Introduction}

A repeated game allows two players to coordinate their answers across
many independently sampled question pairs. Even when their success in a
single game is bounded away from one, the wins in a joint strategy need
not be independent. Threshold parallel repetition asks whether the
probability that the fraction of wins exceeds the single-game value by
any fixed positive amount is exponentially small. Such a bound permits
error amplification even when an intended strategy sometimes loses
individual games.

\subsection{Games and the main theorem}

A finite two-player one-round game consists of nonempty finite question
sets $\mathcal X,\mathcal Y$, nonempty finite answer sets
$\mathcal A,\mathcal B$, a probability distribution $\mu$ on
$\mathcal X\times\mathcal Y$, and a predicate
\[
 V:\mathcal X\times\mathcal Y\times\mathcal A\times\mathcal B
       \longrightarrow\{0,1\}.
\]
The referee samples $(x,y)$ from $\mu$, sends $x$ to Alice and $y$ to
Bob, and accepts answers $(a,b)$ when $V(x,y,a,b)=1$. The players do
not communicate. A finite-dimensional entangled strategy consists of
a unit vector $\psi\in\mathcal H_A\otimes\mathcal H_B$ and, for each
local question, a positive operator-valued measurement (POVM):
positive operators $M_x^a,N_y^b$ satisfying
$\sum_aM_x^a=I$ and $\sum_bN_y^b=I$. Its success probability is
\[
 \sum_{x,y,a,b}\mu(x,y)V(x,y,a,b)
     \langle\psi,M_x^a\otimes N_y^b\psi\rangle.
\]
Using pure states causes no loss, since a finite-dimensional mixed state
can be purified by adding a finite local register. The entangled value
$\omega^*(G)$ is the supremum of this expression over all finite
dimensions and all such strategies; the supremum need not be attained.

In $k$ repetitions the referee samples
$(X_1,Y_1),\ldots,(X_k,Y_k)$ independently from $\mu$ and sends each
player the entire local question string. A strategy may use arbitrary
joint POVMs $M_{\mathbf x}^{\mathbf a}$ and
$N_{\mathbf y}^{\mathbf b}$, with answer strings in
$\mathcal A^k$ and $\mathcal B^k$. Write
\[
 I_i=V(X_i,Y_i,A_i,B_i),\qquad W_k=\sum_{i=1}^k I_i,
 \qquad d=|\mathcal A||\mathcal B|.
\]
All logarithms below are natural.

\begin{theorem}[Threshold parallel repetition]\label{thm:main}
There is a universal constant $\kappa_0>0$ such that the following holds.
Let $G$ be any finite two-player one-round game with Boolean acceptance
predicate and entangled value $v=\omega^*(G)<1$. For every
$0<\delta<1-v$, every integer $k\ge1$, and every finite-dimensional
joint strategy for $k$ repetitions,
\begin{equation}\label{eq:main}
 \Pr\!\left[W_k\ge\lceil(v+\delta)k\rceil\right]
 \le \exp\!\left(-\frac{\kappa_0\delta^5}
                         {1+\log d}\,k\right).
\end{equation}
The same inequality holds for the supremum over those strategies.
\end{theorem}

The distribution $\mu$ may be correlated and may have zero entries or
disconnected support. Neither the question alphabet sizes nor the
strategy dimension enter the displayed rate. We obtain this theorem by
combining a transfer from conditioning bounds to threshold bounds with
Song's conditioning estimate~\cite[Lemma~4.4]{song2026}. The transfer is
proved in Section~\ref{sec:threshold}; its tested-set selection uses
only the joint law of the binary win indicators.

A second bound records the geometry of a fixed question distribution.
Form the graph whose vertices
are the positive-probability pairs $H=\{(x,y):\mu(x,y)>0\}$, joining two
pairs when they share either question. For its connected components
$H_j$, let $\mu_j=\mu(\,\cdot\mid H_j)$ and
$\mu_{j,\min}=\min_{z\in H_j}\mu_j(z)$, and set
\[
 \kappa_\mu=8\max_j\frac{|H_j|-1}{\mu_{j,\min}}.
\]
In particular, a singleton component contributes zero.
Corollary~\ref{cor:support-threshold} proves, under the same hypotheses and for
every $k\ge1$, the complementary estimate
\[
 \Pr\!\left[W_k\ge\lceil(v+\delta)k\rceil\right]
 \le\exp\!\left(-\frac{\delta^3}
 {16384(1+\kappa_\mu)(1+\log d)}\,k\right).
\]
Both estimates apply to the same strategy, so one may use the larger
exponential rate. The cubic bound is proved in
Section~\ref{app:root-rounding} by constructing one reference state per
support component. Its constant depends on the fixed distribution through
$\kappa_\mu$, whereas the fifth-power constant is universal. Both bounds
hold at every repetition length, including when rounding the threshold
to an integer matters. When $v=0$, Lemma~\ref{lem:mean-bound} gives the
sharper probability zero.

\subsection{Related work}

The classical value of a game is its largest success probability using
only shared randomness and local response functions. Raz proved that,
when this value is less than one, the probability of winning every
coordinate decays exponentially under repetition~\cite{raz1998}.
Holenstein's information-theoretic proof introduced correlated sampling
into this argument and improved its quantitative dependence on the gap
from one~\cite{holenstein2009}. Rao then proved exponential concentration
above the classical value for general games~\cite[Theorem~7]{rao2011}.
This threshold conclusion is stronger than an all-wins bound: the event
that every coordinate is won is only a subset of the event that a
specified fraction is won.

For entangled strategies, early repetition results used structure in
the acceptance predicate or in the question distribution. Cleve,
Slofstra, Unger, and Upadhyay proved perfect repetition for XOR games,
whose answers are bits and whose acceptance depends on their parity:
the all-wins value of a product is the product of the individual
values~\cite{csuu2008}. Kempe, Regev, and Toner proved exponential
repetition for unique games, whose constraints are permutations between
the answer sets~\cite{krt2010}. Dinur, Steurer, and Vidick extended
exponential repetition to projection games, in which one player's
answer determines the other player's required answer, and developed
quantum correlated sampling~\cite{dsv2015}. Jain, Pereszl\'enyi, and Yao
treated free games, where the players' questions are independently
distributed~\cite{jpy2014}; Chailloux and Scarpa subsequently improved
the rate~\cite{cs2015}. These results require restrictions on either
the predicate or the question distribution.
For complete-support question distributions, which may be correlated,
Chailloux and Scarpa also proved exponential all-wins decay with a rate
depending on the distribution~\cite[Corollary~1, corrected version]{cs2014}.
That result does not cover arbitrary zero entries in the question law.

Yuen proved polynomial all-wins decay for arbitrary finite two-player
entangled games~\cite[Theorem~1]{yuen2016}. His proof combines
conditioning with classical and quantum correlated sampling; its
tested-coordinate selection already samples indices with
replacement~\cite[Proposition~9, full version]{yuen2016}.
Bavarian, Vidick, and Yuen obtained exponential amplification by
anchoring, a transformation that adds questions to the
game~\cite{bvy2022}. Their threshold theorem concerns these anchored
games~\cite[Section~5.4, Theorem~5.6]{bvy2017}. Rao's and their
concentration arguments are precedents for extracting a threshold bound
from control after conditioning on tested coordinates. Unlike the
anchoring transformation, the present bounds retain the given question
distribution throughout.

Chapter~6 of \emph{Ten Advances in Mathematics and Theoretical Computer
Science} obtains exponential all-wins repetition for arbitrary finite
entangled games~\cite{tenadvances2026}. Its operator entropy estimate
and conditioning construction are the analytic antecedents of the
argument below. We reproduce that chapter's resolvent purification and
its estimate from Definition~4.2 and Lemma~4.3, and adapt its
Section~4 weighted reveal argument to a fixed ordering of the untested
coordinates.
Song replaces the quantum sampler in that framework by a smooth
construction and improves the all-wins gap dependence to a cubic
power~\cite{song2026}. His generic conditioning lemma supplies our
fifth-power threshold theorem. The cubic all-wins result uses an
additional estimate for near-perfect conditional success; it does not
by itself imply a cubic threshold rate above an arbitrary value.

Our contribution is a quantitative transfer from tested-set
conditioning to threshold decay, together with its universal
fifth-power application and the explicit cubic bound for each fixed
question distribution. The universal application uses Song's sampler
independently of the Dinur--Steurer--Vidick sampling theorem, while
retaining the conditioning machinery of Chapter~6. Since $0<\delta<1$,
the weaker rate
$\kappa_0\delta^{13}/(1+\log d)$ already follows from
Theorem~\ref{thm:main}.

\subsection{The proof in outline}

Fix a repeated strategy. For a tested set $C\subsetneq[k]$, let $E_C$
be the event that all coordinates in $C$ are won, and suppose
$p=\Pr(E_C)>0$. Put $m=k-|C|$ and
\[
 \tau=\frac{\log(1/p)+|C|\log d}{m}.
\]
The numerator records the rarity of the conditioning event and the
logarithm of the number of possible tested answer strings. Song's
estimate, stated in Proposition~\ref{prop:song-conditioning}, bounds
the average success on the remaining coordinates by
\begin{equation}\label{eq:outline-conditioning}
 \frac1m\sum_{i\notin C}\Pr(I_i=1\mid E_C)
 \le v+B_4\tau^{1/4}\qquad(0\le\tau\le1),
\end{equation}
where $B_4$ is universal. Thus a small conditioning budget cannot raise
the average win probability far above $v$.

Section~\ref{sec:threshold} first proves the generic transfer. If the
upper tail has too much probability,
sampling a short list of coordinate indices with replacement identifies
a fixed tested set $C$ whose success probability is large and whose
conditional win rate outside $C$ exceeds $v$ by a fixed fraction of
$\delta$. Its budget is small enough to contradict
\eqref{eq:outline-conditioning}. This binary sampling argument gives
the exponential bound for sufficiently large repetition lengths. Taking
independent tensor copies of any shorter strategy then gives the same
rate for every length. The generic transfer turns a conditioning error
of order $\tau^{1/4}$ into the fifth-power rate in the main theorem.

The fixed-distribution cubic bound uses a different construction of a
single-game strategy. Sections~\ref{sec:operators} and
\ref{sec:exposure} represent the conditional answers by measurements on
vectors in common finite-dimensional spaces. Resolvent purifications
preserve the answer probabilities. To compare the vectors, order the
untested coordinates and expose Alice's questions before a chosen
coordinate and Bob's questions after it. Moving this dividing point
gives two entropy telescopes, with total cost at most
$p(\log(1/p)+|C|\log d)$ on each side. The resulting vector has two
nearby approximations, each determined by one player's question and a
transcript of exposed questions and tested answers.

Section~\ref{app:root-rounding} compares these vectors along paths in
each connected component of the question support. The resulting
strategies for the component games can be combined in one fixed
tensor-product state, since each player identifies the component from
the local question. The loss is $O_\mu(\sqrt\tau+\tau)$, and the same
transfer therefore gives cubic dependence on $\delta$.
Corollary~\ref{cor:error-amplification} applies the threshold theorem
to amplify both completeness and soundness.
